\documentclass[11pt]{article}

\usepackage[margin=1in]{geometry}
\usepackage[utf8]{inputenc}
\usepackage[T1]{fontenc}
\usepackage[hidelinks]{hyperref}
\usepackage{booktabs}
\usepackage{array}
\usepackage{amsmath}
\usepackage{microtype}
\usepackage{natbib}
\usepackage{parskip}

\title{Steering ESM2-650M:\\
What Moves, What Does Not, and How to Tell in Advance}
\author{Dmitriy Ivkov}
\date{\today}

\begin{document}
\maketitle

\begin{abstract}
We tried to steer nine sequence properties in the protein language model
ESM2-650M using activation steering, under one fixed method and one acceptance
protocol fixed in advance. One number, measured before any generation---the
amino-acid \emph{compositional separation} between the low- and high-label
groups the steering vector is built from---predicts both whether a property is
steerable and how large the effect is. The proxy's correlation with the true
label does not: the best-validated proxy in the whole set ($r=0.80$, binding)
steered nothing, while the weakest runnable one ($r=0.22$, catalytic turnover)
passed on all three seeds. Effect size grows with separation, from a flat null
to $72\%$ of the natural property gap. Separation is necessary but not
sufficient---two well-separated targets still failed, one because the model
ignores the direction and one because the property is \emph{positional}. That
last case is a structural limit: a vector added uniformly at every position can
change how much of a residue a sequence contains but not where it goes, so this
method can move aggregate compositional properties (transmembrane fraction,
catalytic composition, charge, glycosylation density) but cannot install
positionally-localized features (signal peptides, active-site geometry, domain
boundaries). All results are computational.
\end{abstract}

\section{How we obtained the results}

\paragraph{The steering method (identical for every target).}
The model is ESM2-650M \citep{lin2023esm}. To generate, we mask 30\% of a
sequence's positions and fill them in a single argmax pass
(\texttt{mask\_fill\_generate}). To steer, we build a direction as the
per-layer difference between the mean activations of 150 high-label and 150
low-label proteins, and add it to all 33 layers, the same amount at every
sequence position, scaled by a strength $\alpha$. This is the standard
difference-of-means activation steering used for protein language models
\citep{huang2025steering}. As a control we use a random direction with the same
norm. We evaluate on 150 held-out proteins, score each generation with a
property-specific compositional proxy, and report the effect as a paired
bootstrap contrast (10{,}000 resamples) between the real and random directions.
The masked-fill decoder collapses into repetitive low-complexity output at
$\alpha\ge0.35$, so we sweep dose only on the valid window
$\{0.10,0.15,0.20,0.25\}$; this caps the achievable effect but keeps every
comparison inside a regime where the outputs are real sequences.

\paragraph{The acceptance protocol (fixed before the runs).}
To keep a flattering number from being read as a capability, each target had to
clear six criteria set in advance:
\begin{enumerate}\itemsep2pt
\item beat the random-direction control with a real bootstrap confidence
interval;
\item show a coherent dose-response over at least three strengths;
\item survive \textbf{residue exclusion}---stay significant after the two most
frequently substituted residues are removed from the score (every proxy is a
function of composition, so steering can cheat by just enriching favored
residues);
\item use a proxy validated against real experimental labels \emph{before} the
run ($|r|\ge0.15$, checked in code before the model even loads);
\item beat any prior technique for the property;
\item be adequately powered ($n\ge150$).
\end{enumerate}
Failing 1--4 outright is a \textsc{kill}; passing most but failing a dose check
or residue exclusion is \textsc{ambiguous}; clearing all operative criteria is
a \textsc{pass}. We reran every candidate on three independent seeds.

\paragraph{The separation gate.}
Before spending any GPU time, we measure the mean per-residue amino-acid
composition $L_2$ distance between the low- and high-label groups that build the
vector. It began as a screen for one early null and turned out to be the single
most predictive quantity in the whole effort.

\paragraph{Targets, data, and proxies.}
Nine properties, chosen to span a range of separation and to contrast
properties intrinsic to a single sequence against relational ones. Transmembrane
fraction, DNA-binding, glycosylation density, disulfide content, signal peptide,
zinc-finger and calcium-binding domains came from reviewed UniProt
\citep{uniprot2023}; cross-protein $k_{\mathrm{cat}}$ from DLKcat
\citep{li2022dlkcat}; single-backbone binding from ProteinGym
\citep{notin2023proteingym}. Proxies use published biophysical scales: mean
Kyte--Doolittle hydropathy for membrane character \citep{kyte1982hydropathy},
Lys$+$Arg fraction for nucleic-acid interfaces \citep{luscombe2001dna}, a
glycine-minus-arginine contrast for the turnover--stability tradeoff
\citep{siddiqui2006cold}. Each proxy was validated against the real label on a
held-out split before its run.

\section{Results}

\subsection{Separation predicts steerability and effect; proxy correlation does not}

Table~\ref{tab:scoreboard} lists all nine targets. Sorting by outcome, effect
size follows the pre-run separation and ignores the proxy's label correlation.
The two extremes make the point. Single-backbone binding has the best-validated
proxy anywhere in the project (held-out $r=0.80$, and it generalizes to unseen
positions), yet it steers \emph{nothing} at any strength---because its low- and
high-label sequences are one- and two-residue mutants of the same 188-residue
backbone, so they are nearly identical in composition (separation $0.003$) and
the difference-of-means vector has almost no signal to carry. Catalytic turnover
has the weakest runnable proxy ($r=0.22$) but passes on all three seeds. Proxy
strength and steerability are not even weakly related.

\begin{table}[htbp]
\centering
\caption{Nine steering targets on ESM2-650M, one fixed method, sorted by
outcome. \emph{Sep.}\ is the pre-run compositional distance between the vector
groups; \emph{proxy $r$} is the held-out correlation of the scoring rule with
the real label; \emph{effect} is, where a natural low$\to$high proxy gap is
defined, the fraction of that gap reached at $\alpha=0.25$. Effect tracks
separation, not proxy $r$.}
\label{tab:scoreboard}
\begin{tabular}{@{}l r r l l@{}}
\toprule
Target & Sep. & Proxy $r$ & Outcome & Effect / why it failed \\
\midrule
Binding, single-backbone            & 0.003 & 0.80 & \textsc{kill} & flat null --- no separation \\
Zinc-finger domain                  & 0.028 & 0.29 & \textsc{kill} & direction anti-correlated \\
Calcium-binding (EF-hand)           & 0.038 & 0.16 & \textsc{kill} & model ignores the direction \\
Signal peptide                      & 0.035 & 0.74 & \textsc{kill} & positional, not compositional \\
\addlinespace
Disulfide (Cys fraction)            & 0.026 & 0.17 & pass (2/3) & tiny ($\sim$0.002) \\
Glycosylation density               & 0.039 & 0.41 & pass (3/3) & tiny ($\sim$0.003) \\
DNA-binding propensity              & 0.039 & 0.25 & pass (2/3) & 25\% of natural gap \\
Catalytic turnover ($k_{\mathrm{cat}}$) & 0.02--0.05 & 0.22 & pass (3/3) & moderate \\
Transmembrane fraction              & 0.086 & 0.79 & pass (3/3) & 72\% of natural gap \\
\bottomrule
\end{tabular}
\end{table}

The passing targets calibrate the relationship: separation $0.003 \to$ null,
$0.039 \to$ a quarter of the natural binder gap, $0.086 \to$ nearly
three-quarters of the way to the membrane-protein distribution (proxy
$-0.363 \to +0.266$ at $\alpha=0.25$ against a natural gap of $0.871$). The
binding null was originally read as ``binding is not steerable.'' Rebuilding the
same target from \emph{cross-protein} data---DNA-binding, where a protein binds
the backbone by its own sequence---lifts separation about $12\times$ ($0.003
\to 0.039$) and turns the null into a significant, residue-robust,
dose-responsive effect. So the null was an artifact of single-backbone data, not
a property of binding, and we predicted the recovery from separation before
running it.

\subsection{Separation is necessary but not sufficient}

Screening by separation kills the null mode but does not guarantee success. Two
targets with separation in the passing band still failed, for different reasons.
Calcium-binding (separation $0.038$) produces no significant steering at any
safe strength---a well-separated group whose direction the model's residual
stream simply does not act on. Zinc-finger's cysteine/histidine direction turns
\emph{anti-correlated} with its own proxy after steering. High separation is a
prerequisite; the model must also respond to the direction, and the direction
must encode the property rather than a correlate.

\subsection{The method steers composition, not position}

The sharpest failure is the signal peptide. Its proxy---mean hydropathy over the
first 30 residues---is the best-validated candidate in the whole set (held-out
$r=0.742$), and it gives a clean dose-response (N-terminal hydropathy
$0.059 \to 0.100 \to 0.152 \to 0.251$). It still \textsc{kill}s on residue
exclusion: remove the two dominant substituted residues and the effect
collapses to non-significant. The vector raises N-terminal hydropathy only by
dumping bulk hydrophobic residues across the \emph{entire} sequence; it cannot
\emph{concentrate} them at the N-terminus, which is what makes a signal peptide.
Glycosylation density passes for the mirror-image reason: its proxy is a sequon
count summed over the whole sequence, a genuinely aggregate quantity a
mean-pooled vector can nudge.

The general statement: a direction that is averaged over positions and added
identically to every position can change \emph{how much} of a residue type a
sequence has, but not \emph{where} it sits. Aggregate compositional
properties---transmembrane fraction, catalytic composition, DNA-binding charge,
glycosylation density---are reachable. Positionally-localized features---signal
peptides, and by the same argument active-site geometry, domain boundaries,
termini motifs---are not, for any proxy, until the injection itself is made
position-aware.

\subsection{What the acceptance protocol caught elsewhere}

The predictor is only trustworthy because the same protocol threw out flattering
numbers in the studies that fed it, and each rejection changed a conclusion:
\begin{itemize}\itemsep2pt
\item An immune-endpoint proxy correlated $r=0.427$ with an allele-aggregated
binding label but only $0.100$ with measured T-cell response, and flipped from
$+0.379$ to $-0.323$ when validation folds were grouped by source organism ---
it measured the wrong endpoint, and was stopped before generation.
\item A layer-subset thermostability run looked competitive only because the
all-layer arm's outputs had collapsed under the low-complexity filter; inside
the shared valid range the subset kept just $43$--$59\%$ of the full effect.
\item A solubility contrast of $+0.0125$ fell to $+0.00035$ (interval crossing
zero) once its two dominant substituted residues were excluded.
\item A disorder effect repeated across three runs but survived residue
exclusion in only two of them.
\end{itemize}
These are the same guards---endpoint validity, evaluable denominator,
composition dependence, whole-run stability---that make a passing row in
Table~\ref{tab:scoreboard} mean something.

\section{Takeaways and limitations}

For this method, steerability and its ceiling are predictable before any compute
from the compositional separation of the vector-building groups, and an entire
class of properties (positional features) is ruled out on structural grounds,
not by trying and failing one at a time. The next moves are to change the
intervention---position-aware injection---rather than the target list, and to
check the strongest compositional gain (transmembrane) against a real assay.

The limits are firm. Every endpoint is computational: no wet-lab or structural
assay shows that any generated sequence actually has more thermostability,
turnover, membrane character, binding, or glycosylation; every ``effect'' is a
change in a compositional proxy. Everything uses one model, one decoder, and one
injection geometry, so the composition-only limit is specific to that geometry.
The usable dose range is narrow because the decoder degenerates above
$\alpha=0.25$. The proxies are compositional and non-mechanistic---a pass shows
the model can be pushed toward an amino-acid signature that correlates with the
property, not that structure or function is preserved. The separation predictor
is calibrated on nine targets from one model; whether the same ordering holds
for other models or injection schemes is untested.

\bibliographystyle{plainnat}
\bibliography{reference}

\end{document}
